\section{Methodology}\label{sec:method}
% NGUOI VIET: LR (Hai) + MS (Phuc)
% Ban nhap tu RW (Loc), tong hop tu proposal.md Sec5, amendment v1.1/v1.2, va notes.md.
% Can Hai/Phuc doc lai + xac nhan chi tiet trien khai truoc khi nop.

\subsection{Dataset}
The dataset consists of \textbf{120 real-world functions} (60 Java, 60 Python) with
cyclomatic complexity $5 \le \mathrm{CC} \le 10$, measured with Lizard. The original
proposal specified Defects4J and CodeXGLUE as data sources; during dataset preparation the
team found that CodeXGLUE is a code-intelligence benchmark (summarisation/translation) and
not a runnable function corpus suited to coverage/mutation measurement, and that the
surveyed literature does not consistently use this pairing. This was recorded as
amendment~v1.2 (2026-07-02, approved before the full run) and resolved by mining functions
directly from ten pinned open-source repositories at fixed commits instead: eight Java
projects overlapping with Defects4J's subject programs (\texttt{commons-cli},
\texttt{commons-math}, \texttt{commons-csv}, \texttt{commons-collections}, \texttt{gson},
\texttt{jsoup}, \texttt{joda-time}, \texttt{jfreechart}) and two Python projects
(\texttt{flask}, \texttt{requests}). Each repository is cloned at a pinned commit with full
provenance (URL, licence, commit hash, download date) recorded separately, so the dataset is
exactly reproducible. The Python split is uneven by design of the source ecosystem (more
functions come from \texttt{flask} than \texttt{requests}); this is disclosed as a threat to
external validity (Section~\ref{sec:threats}) rather than corrected by resampling, to avoid
altering the dataset after any measurement had begun.

\subsection{Test generation}
\textbf{LLM.} We evaluate \texttt{gpt-4o-mini-2024-07-18} (pinned snapshot),
\texttt{temperature=0}, \texttt{top\_p=1}, \texttt{max\_tokens=2048}, one-shot prompting (a
single worked exemplar precedes the target function in the prompt). The original proposal
specified \texttt{gpt-4o}; the team's OpenAI account only had access to \texttt{gpt-4o-mini}
at execution time, an objective access constraint recorded as amendment~v1.1
(2026-06-28, approved before the full run) rather than a post-hoc substitution --- a
legitimate academic precedent exists for evaluating \texttt{gpt-4o-mini} specifically (one
of the papers in our own related-work table does so). All research questions, metrics,
thresholds, and statistical tests were kept unchanged by both amendments.

\textbf{Baselines.} Java: EvoSuite and Randoop, both under an equal 60-second
per-class/per-function search budget (an initial 60s-vs-10s asymmetry between EvoSuite and
Randoop was identified during pilot testing and corrected before the full run;
Section~\ref{sec:threats}). Python: Pynguin (DYNAMOSA algorithm, 90-second search budget per
module); Hypothesis was listed as a candidate baseline in the original proposal but was not
exercised in the reported results.

\subsection{Measurement}
Generated tests are executed \emph{inside the original project}, not in isolation: each
function is measured by building/running the actual repository at its pinned commit,
because a substantial fraction of functions reference names and behaviour that only resolve
correctly inside their real module (Section~\ref{sec:threats} documents the pilot finding
that motivated this choice). For Java, tests are compiled and run with Maven; branch
coverage is collected with JaCoCo and mutation testing with PIT, both attributed \emph{per
function} by intersecting the tool's line-level report with each function's
\texttt{[start\_line, end\_line]} range in the pinned source file --- necessary because
JaCoCo/PIT report at class granularity while our unit of analysis is the function. Because
GPT-generated Java suites have a much higher per-file compile-failure rate than
tool-generated ones, each GPT-authored test is compiled and measured \emph{individually}
(one file per function per run), so that a single compile error cannot be mistaken for a
failure in a co-located, otherwise-valid test in the same class. For Python, functions are
measured inside the installed module (editable install of the pinned \texttt{flask}/
\texttt{requests} checkout); branch coverage uses \texttt{coverage.py} and mutation testing
uses a lightweight AST-level mutation operator (binary-operator, comparison-operator, and
boolean-operator swaps, boolean/numeric-constant perturbation) restricted to statements
within each function's line range, applied identically to both the GPT-4o-mini and Pynguin
arms so the two are measured with the same instrument.

A test suite is scored \textsc{invalid} (branch coverage $=0$, mutation score $=0$) if it
fails to compile/import or fails to execute (proposal~\S5.1). A stronger, mandatory
\emph{green-on-original} check is additionally enforced before any mutation analysis: a
suite that does not pass on the unmutated source is rejected outright, since scoring
mutation kills for an already-failing suite would silently inflate the mutation score
(Section~\ref{sec:threats} documents a concrete case this check caught). We further
distinguish, at the reporting level, between a suite that is \textsc{invalid} and a suite
that \emph{compiles and runs but exercises none of the target function's lines}
(branch coverage $=0$ despite \texttt{compiled=1}); both score 0\% on our metrics, but they
are different failure mechanisms and are reported separately in Section~\ref{sec:results}.

\subsection{Statistical analysis}
All tests use $\alpha = 0.05$ and non-parametric methods, chosen because per-function
coverage/mutation distributions are not assumed normal and are bounded in $[0,100]$.
\textbf{RQ1} (coverage $\ge 80\%$): one-sample Wilcoxon signed-rank test of branch coverage
against 80. \textbf{RQ2} (mutation $\ge 60\%$ and vs.\ baseline): a one-sample Wilcoxon
against 60 for the threshold sub-question, and a paired Wilcoxon signed-rank test
(gpt-4o-mini vs.\ each baseline, matched by \texttt{function\_id}) for the comparative
sub-question, with matched-pairs rank-biserial correlation $r$ as effect size.
\textbf{RQ3} (complexity vs.\ quality): Spearman rank correlation between CC and each
quality metric. Every statistic is reported at two levels: \emph{all} ($N=60$ per language;
\textsc{invalid}/non-touching functions scored 0, per the pre-registered rule) and
\emph{effective} (the conditional subset of functions whose suite compiled and touched the
target), since the effective-level figures alone would overstate raw capability
(Section~\ref{sec:discussion}).
